% Solutions/hints for ch17 exercises (assembled into Appendix C)

\begin{solution}{exr:ch17-hand}
The nearest vertex is $\vect{v}_8=(4.4,3.0)$ at distance $0.721<\eta$, so $\state_{\mathrm{new}}=\state_{\mathrm{rand}}=(5.0,2.6)$, and the segment $\vect{v}_8\state_{\mathrm{new}}$ is free. $\Near$ with $r_n=2$ returns $\set{\vect{v}_4,\vect{v}_7,\vect{v}_8,\vect{v}_9}$ at distances $1.393$, $1.772$, $0.721$ and $1.217$. \textsc{ChooseParent} sorts the candidates by $\mathrm{Cost}(\vect{v})+\norm{\state_{\mathrm{new}}-\vect{v}}$: $\vect{v}_7$ ($4.043+1.772=5.815$), $\vect{v}_4$ ($6.141$), $\vect{v}_8$ ($7.173$), $\vect{v}_9$ ($8.799$). The first candidate $\vect{v}_7$ has a free segment, so it is chosen after a single collision check and $\state_{\mathrm{new}}$ becomes $\vect{v}_{10}$ with cost $5.815$; the nearest vertex $\vect{v}_8$ would have given $7.173$. \textsc{Rewire} tests the others through $\vect{v}_{10}$: $\vect{v}_4$ would cost $5.815+1.393=7.208>4.748$ and $\vect{v}_8$ would cost $6.536>6.451$, both rejected by the improvement test without a collision check; $\vect{v}_9$ would cost $5.815+1.217=7.032<7.583$ and its segment is free, so its parent changes from $\vect{v}_8$ to $\vect{v}_{10}$ and its cost drops to $7.032$. It has no children, so nothing else changes. \code{tiny\_example(x\_rand=(5.0, 2.6))} returns exactly these values.
\end{solution}

\begin{solution}{exr:ch17-radius}
(a) The cap stops being active when $13.3\sqrt{\log n/n}=1$, that is at
$n=1\,264$ vertices. At $n=10^4$ the expected size of the \textsc{Near} set
is $\zeta_2\gamma^2\log n/\mu(\Xfree)=\pi\cdot13.3^2\cdot\log(10^4)/76=67$
vertices. (b) With $d=3$ and $\zeta_3=4\pi/3$,
$\gamma^*_{\rrtstar}=2(4/3)^{1/3}(716/\zeta_3)^{1/3}=12.22$ for the true free
volume and $13.66$ for the box volume $\mu(\mathcal{X})=1\,000$, so the
$\gamma=15$ of the code is admissible for either. The cap $\eta=1.5$ stops
being active when $15(\log n/n)^{1/3}=1.5$, that is $\log n/n=10^{-3}$ and
$n\approx9\,119$ vertices. (c) A ball of radius $r_n=\gamma\log n/n$ has area
$\zeta_2\gamma^2(\log n/n)^2$, so the expected number of vertices in it is
$n\cdot\zeta_2\gamma^2(\log n/n)^2/\mu(\Xfree)=\zeta_2\gamma^2(\log n)^2/(n\,\mu(\Xfree))\to0$:
the \textsc{Near} set is empty from some $n$ on, \rrtstar degenerates to
\rrt and stops improving. The exponent $1/d$ is what keeps the count
growing like $\log n$.
\end{solution}

\begin{solution}{exr:ch17-ellipsoid}
(a) Put the foci at $\vect{f}_\pm=(\pm c_{\min}/2,0,\dots,0)$ and write $\vect{y}=(y_1,\vect{y}_\perp)$. The condition $\norm{\vect{y}-\vect{f}_-}+\norm{\vect{y}-\vect{f}_+}\le c_{\mathrm{best}}$ depends on $\vect{y}$ only through $y_1$ and $\norm{\vect{y}_\perp}$, so the set is a body of revolution about the first axis, and in the plane spanned by $\vect{e}_1$ and any perpendicular direction it is the classical ellipse with foci at distance $c_{\min}/2$ from the centre and string length $c_{\mathrm{best}}$: squaring twice gives $y_1^2/a^2+\norm{\vect{y}_\perp}^2/b^2\le1$ with $a=c_{\mathrm{best}}/2$ and $b^2=a^2-c_{\min}^2/4$, that is $b=\tfrac12\sqrt{c_{\mathrm{best}}^2-c_{\min}^2}$. This is the image of the unit ball under $\mat{L}=\diag(a,b,\dots,b)$. The world frame has the focal axis along $\vect{a}_1$ and the centre at $\state_{\mathrm{centre}}$, so $\state=\mat{C}\vect{y}+\state_{\mathrm{centre}}$ with any rotation $\mat{C}\vect{e}_1=\vect{a}_1$; rotations preserve distances, so the foci land on $\state_{\mathrm{init}}$ and $\state_{\mathrm{goal}}$.
(b) If $\vect{X}$ has density $p$ on $B$ and $\vect{Y}=\mat{A}\vect{X}+\vect{t}$, the change-of-variables formula gives $\vect{Y}$ the density $p(\mat{A}\inv(\vect{y}-\vect{t}))/\abs{\det\mat{A}}$ on $\mat{A}B+\vect{t}$. For $p\equiv1/\mu(B)$ this is the constant $1/(\mu(B)\abs{\det\mat{A}})=1/\mu(\mat{A}B+\vect{t})$: uniform. With $\mat{A}=\mat{C}\mat{L}$, $\abs{\det\mat{A}}=ab^{d-1}$.
(c) For $\vect{X}$ uniform in the unit ball, $\Prob(\norm{\vect{X}}\le\rho)=\rho^d$ (volume ratio) and the direction is independent of the radius and uniform on the sphere. With $w\sim U(0,1)$, $\Prob(w^{1/d}\le\rho)=\Prob(w\le\rho^d)=\rho^d$, and $\vect{u}/\norm{\vect{u}}$ is uniform on the sphere because the standard Gaussian density $\propto e^{-\norm{\vect{u}}^2/2}$ is invariant under rotations; the two factors are independent, so the product is uniform in the ball.
(d) $\mat{C}=\mat{U}\mat{D}\mat{V}\T$ is a product of orthogonal matrices, hence orthogonal, and $\det\mat{C}=\det\mat{U}\cdot(\det\mat{U}\det\mat{V})\cdot\det\mat{V}=(\det\mat{U})^2(\det\mat{V})^2=1$. The matrix $\vect{a}_1\vect{e}_1\T$ has rank one, so its SVD has $\sigma_1=1$, $\vect{u}_1=s\vect{a}_1$ and $\vect{v}_1=s\vect{e}_1$ for a common sign $s$, and the remaining columns of $\mat{V}$ are orthogonal to $\vect{e}_1$; thus $\mat{V}\T\vect{e}_1=s\vect{e}_1$, $\mat{D}$ leaves $\vect{e}_1$ unchanged ($d\ge2$), and $\mat{C}\vect{e}_1=s\mat{U}\vect{e}_1=s\vect{u}_1=s^2\vect{a}_1=\vect{a}_1$. In 2D the only rotations are $\bigl[\begin{smallmatrix}\cos\phi&-\sin\phi\\\sin\phi&\cos\phi\end{smallmatrix}\bigr]$, and the first column must be $\vect{a}_1=(\cos\phi,\sin\phi)$; the code returns $\bigl[\begin{smallmatrix}0.707&-0.707\\0.707&0.707\end{smallmatrix}\bigr]$ for the diagonal start--goal pair of the worked example.
\end{solution}

\begin{solution}{exr:ch17-invariant}
\emph{No cycles.} Suppose \textsc{Rewire} considers making $\state_{\mathrm{new}}$ the parent of $\vect{v}$ while $\state_{\mathrm{new}}$ is a descendant of $\vect{v}$, that is, the tree path from $\state_{\mathrm{init}}$ to $\state_{\mathrm{new}}$ passes through $\vect{v}$. Costs add along tree paths, so $\mathrm{Cost}(\state_{\mathrm{new}})=\mathrm{Cost}(\vect{v})+\ell$ where $\ell$ is the length of the tree path from $\vect{v}$ to $\state_{\mathrm{new}}$, and $\ell\ge\norm{\state_{\mathrm{new}}-\vect{v}}$ by the triangle inequality. The candidate cost is therefore $\mathrm{Cost}(\state_{\mathrm{new}})+\norm{\vect{v}-\state_{\mathrm{new}}}\ge\mathrm{Cost}(\vect{v})+2\norm{\vect{v}-\state_{\mathrm{new}}}\ge\mathrm{Cost}(\vect{v})$, and the strict test $c<\mathrm{Cost}(\vect{v})$ fails. \emph{Example.} Root $\vect{v}_0=(0,0)$ with child $\vect{a}=(1,0)$ and grandchild $\vect{p}=(2,0)$, so $\mathrm{Cost}(\vect{a})=1$, $\mathrm{Cost}(\vect{p})=2$; obstacle $[1.2,1.4]\times[0.3,0.5]$; sample $\state_{\mathrm{new}}=(1.6,0.8)$; $r_n=1$, so $\Near=\set{\vect{a},\vect{p}}$ (distances $1.0$ and $0.894$). \textsc{ChooseParent} tests $\vect{a}$ first ($1+1.0=2.0$) but the segment crosses the obstacle; $\vect{p}$ ($2+0.894=2.894$) is free, so $\state_{\mathrm{new}}$ hangs under $\vect{p}$ with cost $2.894$. \textsc{Rewire} then considers $\vect{a}$: the cost via $\state_{\mathrm{new}}$ would be $2.894+1.0=3.894>1$, so the candidate that would have closed the cycle $\vect{a}\to\vect{p}\to\state_{\mathrm{new}}\to\vect{a}$ is rejected before any collision check.
\emph{Order.} Let $\vect{v}$ be an ancestor of $\vect{u}$ and both be candidates with $c_{\vect{v}}=\mathrm{Cost}(\state_{\mathrm{new}})+\norm{\vect{v}-\state_{\mathrm{new}}}<\mathrm{Cost}(\vect{v})$ and $c_{\vect{u}}=\mathrm{Cost}(\state_{\mathrm{new}})+\norm{\vect{u}-\state_{\mathrm{new}}}<\mathrm{Cost}(\vect{u})$. If $\vect{u}$ is processed first, both are rewired and both end up as children of $\state_{\mathrm{new}}$. If $\vect{v}$ is processed first, $\mathrm{Cost}(\vect{u})$ drops to $c_{\vect{v}}+\ell(\vect{v},\vect{u})$ where $\ell(\vect{v},\vect{u})\ge\norm{\vect{u}-\vect{v}}$ is the tree distance; but $c_{\vect{u}}\le\mathrm{Cost}(\state_{\mathrm{new}})+\norm{\vect{v}-\state_{\mathrm{new}}}+\norm{\vect{u}-\vect{v}}\le c_{\vect{v}}+\ell(\vect{v},\vect{u})$, with equality only when $\state_{\mathrm{new}}$, $\vect{v}$, $\vect{u}$ are collinear in that order, so $\vect{u}$ is still rewired and the final tree is the same. Candidates that are not rewired are unaffected by the others, because their costs do not change. The work differs: processing $\vect{v}$ first propagates through $\vect{u}$'s subtree twice. Processing candidates in decreasing order of $\mathrm{Cost}$ (descendants before ancestors) touches every subtree once.
\end{solution}

\begin{solution}{exr:ch17-fraction}
(a) With $c_{\min}=11.31$ and \cref{eq:ch17-axes}, the ellipse of the $30\times30$ field has area $\pi\cdot\frac{c_{\mathrm{best}}}{2}\cdot\frac{\sqrt{c_{\mathrm{best}}^2-c_{\min}^2}}{2}$: $410$ at $c_{\mathrm{best}}=24.3$ ($45.6\,\%$ of $900$), $231$ at $19.1$ ($25.7\,\%$) and $91$ at $14.0$ ($10.1\,\%$). Rejection sampling would need $1/0.456=2.2$, $3.9$ and $9.9$ uniform draws per accepted sample. (b) In 3D the volume is $\frac{4\pi}{3}r_1r_2^2$ with $r_1=c_{\mathrm{best}}/2$ and $r_2=\frac12\sqrt{c_{\mathrm{best}}^2-15.59^2}$: $3\,768$ at $23.45$ and $2\,217$ at $21.07$, larger than the box of volume $1\,000$, so the informed set is (nearly) the whole box---Monte Carlo with $2\times10^6$ points in $[0,10]^3$ puts $100\,\%$ and $98.4\,\%$ of it inside, that is $1.00$ and $1.02$ draws per accepted sample---and the code samples the box and rejects by $\hat f$; $410$ at $17.0$ and $109$ at $16.0$, that is $41\,\%$ and $10.9\,\%$ of the box volume. These are volumes of the ellipsoid; the ellipsoid of $c_{\mathrm{best}}=17.0$ still pokes out of the box, so only $38.7\,\%$ of the box lies inside the informed set and rejection sampling of the box needs $2.6$ draws per accepted sample, while the ellipsoid of $16.0$ fits inside and needs $9.2$. The conjugate axes enter with the power $d-1$, which is why the fraction falls faster in 3D. (c) The map is inside the ellipse as long as its farthest corners are, and the corners $(10,0)$ and $(0,10)$ have $\hat f=2\sqrt{81+1}=18.11$ (the corners $(0,0)$ and $(10,10)$ have $\hat f=14.14$). Below $c_{\mathrm{best}}=18.11$ the corners drop out, but the excluded area is tiny; at $c_{\mathrm{best}}=17.32$ the code rejected $8$ of $2\,657$ draws.
\end{solution}
